\chapter{Minimi quadrati e stima ricorsiva}

Dimentichiamoci ora della probabilità. In questo capitolo i vettori contano per quello che sono, e non perché sono
considerati realizzazioni di vettori aleatori: è il punto di vista della \textbf{statistica descrittiva}. Si hanno dei
dati misurati e si cerca il modello lineare che li descrive meglio, nel senso dei \textbf{minimi quadrati}. È lo stesso
problema dello stimatore lineare MMSE, con le medie sui dati al posto delle medie statistiche. Alla fine vedremo come
aggiornare la stima quando arriva un dato nuovo, senza rifare tutti i conti da capo.

% ============================================================
\section{Il modello lineare e la stima ai minimi quadrati}
% ============================================================

Un \term{insieme di addestramento} (\emph{training set}) è una collezione di $p$ vettori $n$-dimensionali $\vett{x}(1),
\ldots, \vett{x}(p)$, che si possono organizzare come colonne di una matrice $n \times p$:
\[ \matr{X} = \begin{bmatrix} x_1(1) & x_1(2) & \cdots & x_1(p) \\ x_2(1) & x_2(2) & \cdots & x_2(p) \\ \vdots & \vdots
   & \ddots & \vdots \\ x_n(1) & x_n(2) & \cdots & x_n(p) \end{bmatrix} = \left[ \vett{x}(1)\ \ \vett{x}(2)\ \ \cdots\ \
   \vett{x}(p) \right] \in \R^{n \times p} . \]
Si conoscono anche $p$ valori misurati del parametro $\theta$, uno per ciascun vettore dell'insieme di addestramento:
$\vett{y} = [\theta(1), \theta(2), \ldots, \theta(p)]^T \in \R^p$. Per esempio, ogni $\vett{x}(i)$ può raccogliere $n$
caratteristiche di una casa (superficie, numero di stanze, piano) e $\theta(i)$ il suo prezzo.

Si vuole adattare ai dati un \textbf{modello lineare}
\[ \theta(i) = \vett{a}^T \vett{x}(i) + \epsilon_i , \qquad i = 1, \ldots, p , \]
dove $\vett{a}$ è un vettore $n$-dimensionale di coefficienti da scegliere ed $\epsilon_i$ è l'errore del modello sul
dato $i$. Si sceglie $\vett{a}$ in modo da minimizzare la somma dei quadrati degli errori,
\[ \norm{\vett{\epsilon}}^2 = \sum_{i=1}^{p} \left( \vett{a}^T \vett{x}(i) - \theta(i) \right)^2 , \qquad \vett{\epsilon}
   = [\epsilon_1, \ldots, \epsilon_p]^T , \]
dove $\norm{\vett{v}}^2 = \vett{v}^T \vett{v} = \sum_i v_i^2$ è il quadrato della lunghezza del vettore (Capitolo~16).

La somma si scrive in modo compatto con la matrice dei dati. Il prodotto $\matr{X}^T \vett{a}$ ($p \times n$ per $n
\times 1$) è un vettore $p \times 1$: la riga $i$ di $\matr{X}^T$ è $\vett{x}(i)^T$, quindi il suo elemento $i$-esimo è
$\vett{x}(i)^T \vett{a} = \vett{a}^T \vett{x}(i)$ (un numero è uguale al suo trasposto). Allora
\[ \matr{X}^T \vett{a} - \vett{y} = \begin{bmatrix} \vett{x}(1)^T \vett{a} - \theta(1) \\ \vdots \\ \vett{x}(p)^T
   \vett{a} - \theta(p) \end{bmatrix} \qquad\text{e}\qquad \sum_{i=1}^{p} \left( \vett{a}^T \vett{x}(i) - \theta(i)
   \right)^2 = \norm{\matr{X}^T \vett{a} - \vett{y}}^2 . \]

\teorema[minimiquadrati]{La stima ai minimi quadrati}{
  Se la matrice simmetrica $n \times n$ $\matr{X}\matr{X}^T$ è invertibile, il vettore che minimizza
  $\norm{\matr{X}^T \vett{a} - \vett{y}}^2$ è
  \[ \vett{a}_{\text{LS}} = \left( \matr{X} \matr{X}^T \right)^{-1} \matr{X}\, \vett{y} , \]
  la \term{stima ai minimi quadrati} (LS, \emph{Least Squares}). È la soluzione delle \term{equazioni normali}
  $\matr{X}\matr{X}^T \vett{a} = \matr{X}\vett{y}$.
}

\begin{dimostrazione}
\passo{Il quadrato di una differenza} Per due vettori $\norm{\vett{\alpha} - \vett{\beta}}^2 = (\vett{\alpha} -
\vett{\beta})^T (\vett{\alpha} - \vett{\beta}) = \norm{\vett{\alpha}}^2 + \norm{\vett{\beta}}^2 - 2\vett{\alpha}^T
\vett{\beta}$, come $(a - b)^2 = a^2 + b^2 - 2ab$. Con $\vett{\alpha} = \matr{X}^T \vett{a}$ e $\vett{\beta} = \vett{y}$,
e il trasposto di un prodotto $(\matr{X}^T \vett{a})^T = \vett{a}^T \matr{X}$:
\[ \norm{\matr{X}^T \vett{a} - \vett{y}}^2 = \vett{a}^T \matr{X} \matr{X}^T \vett{a} + \norm{\vett{y}}^2 - 2\vett{a}^T
   \matr{X} \vett{y} . \]
\passo{Il gradiente} $\matr{X}\matr{X}^T$ è simmetrica (è uguale alla sua trasposta). Con le formule del Capitolo~23
(il secondo termine non dipende da $\vett{a}$):
\[ \nabla_{\vett{a}} \norm{\matr{X}^T \vett{a} - \vett{y}}^2 = 2\matr{X}\matr{X}^T \vett{a} - 2\matr{X}\vett{y} =
   \vett{0} \quad\Longrightarrow\quad \vett{a} = \left( \matr{X}\matr{X}^T \right)^{-1} \matr{X}\, \vett{y} . \]
\passo{È un minimo} La somma dei quadrati è una funzione di secondo grado a forma di scodella: $\vett{q}^T \matr{X}
\matr{X}^T \vett{q} = \norm{\matr{X}^T \vett{q}}^2 \geq 0$ per ogni $\vett{q}$, quindi la matrice è definita non
negativa (Capitolo~16), e l'unico punto a gradiente nullo è il fondo della scodella.
\end{dimostrazione}

La formula richiede che la matrice $\matr{X}\matr{X}^T$ sia invertibile. In sostanza servono $p \geq n$ dati (almeno
tante equazioni quante sono le incognite) e che le $n$ colonne della matrice $\matr{X}^T$, di dimensioni $p \times n$,
siano \term{linearmente indipendenti}: nessuna caratteristica deve essere una combinazione lineare delle altre (se due
caratteristiche dicessero la stessa cosa, per esempio la superficie in metri quadri e in centimetri quadri, non si
saprebbe come dividere il peso fra loro).

\info{Il legame con lo stimatore lineare MMSE}{
  Dividendo per $p$, $\frac{1}{p}\matr{X}\matr{X}^T = \frac{1}{p}\sum_i \vett{x}(i)\vett{x}(i)^T$ e $\frac{1}{p}\matr{X}
  \vett{y} = \frac{1}{p}\sum_i \theta(i)\, \vett{x}(i)$ sono le medie campione di $\vett{X}\vett{X}^T$ e di
  $\Theta\vett{X}$: per la legge dei grandi numeri tendono a $\matr{R} = E[\vett{X}\vett{X}^T]$ e a $E[\Theta\vett{X}]$.
  La stima LS è la versione «con i dati» dello stimatore lineare MMSE del Capitolo~23 (senza la costante $b$; per averla
  basta aggiungere a ogni $\vett{x}(i)$ una componente sempre uguale a 1, come nell'esempio che segue).
}

\begin{esempio}{La retta che passa più vicina a sei punti}
Si cercano i coefficienti della retta $\theta = a_1 + a_2 t$ che approssima meglio sei punti misurati:
\begin{center}
\renewcommand{\arraystretch}{1.25}
\begin{tabular}{@{}c|cccccc@{}}
  $t_i$ & 1 & 2 & 3 & 4 & 5 & 6 \\ \hline
  $\theta_i$ & $2{,}1$ & $3{,}9$ & $6{,}2$ & $7{,}8$ & $10{,}1$ & $12{,}0$
\end{tabular}
\end{center}
È il modello lineare con $\vett{x}(i) = [1, t_i]^T$ (la componente costante 1 moltiplica $a_1$). Quindi
\[ \matr{X}\matr{X}^T = \sum_i \vett{x}(i) \vett{x}(i)^T = \begin{bmatrix} p & \sum t_i \\ \sum t_i & \sum t_i^2
   \end{bmatrix} = \begin{bmatrix} 6 & 21 \\ 21 & 91 \end{bmatrix} , \qquad \matr{X}\vett{y} = \begin{bmatrix} \sum
   \theta_i \\ \sum t_i \theta_i \end{bmatrix} = \begin{bmatrix} 42{,}1 \\ 182{,}2 \end{bmatrix} . \]
Con la formula dell'inversa $2 \times 2$ (determinante $6 \cdot 91 - 21^2 = 105$):
\[ \vett{a}_{\text{LS}} = \frac{1}{105} \begin{bmatrix} 91 & -21 \\ -21 & 6 \end{bmatrix} \begin{bmatrix} 42{,}1 \\
   182{,}2 \end{bmatrix} = \frac{1}{105} \begin{bmatrix} 3831{,}1 - 3826{,}2 \\ -884{,}1 + 1093{,}2 \end{bmatrix} =
   \begin{bmatrix} 0{,}047 \\ 1{,}991 \end{bmatrix} . \]
La retta è $\theta \approx 0{,}047 + 1{,}991\, t$. Gli errori $\epsilon_i$ sono i segmenti verticali della figura; la
somma dei loro quadrati, $0{,}107$, è la più piccola possibile fra tutte le rette.
\begin{center}
\begin{tikzpicture}
\begin{axis}[width=10.5cm, height=6cm, xmin=0, xmax=7, ymin=0, ymax=13.5, xtick={1,2,3,4,5,6}, ytick={2,4,6,8,10,12},
  xlabel=$t$, ylabel=$\theta$, axis lines=left, tick label style={font=\scriptsize}, label style={font=\small},
  grid=major, grid style={mspigray!20}]
  \addplot[mspiB, line width=1.1pt, domain=0.3:6.7] {0.0467 + 1.9914*x};
  \foreach \t/\y in {1/2.1, 2/3.9, 3/6.2, 4/7.8, 5/10.1, 6/12.0} {
    \edef\temp{\noexpand\draw[mspiC, line width=1pt] (axis cs:\t,\y) -- (axis cs:\t,{0.0467 + 1.9914*\t});}\temp
  }
  \addplot[only marks, mark=*, mark size=2pt, mspiA] coordinates {(1,2.1) (2,3.9) (3,6.2) (4,7.8) (5,10.1) (6,12.0)};
  \node[font=\small, mspiB!80!black, anchor=west] at (axis cs:4.3,6.3) {$\theta = 0{,}047 + 1{,}991\, t$};
\end{axis}
\end{tikzpicture}
\end{center}
(Gli errori sono così piccoli che nella figura i segmenti verdi si vedono appena: i punti sono quasi allineati.)
\end{esempio}

% ============================================================
\section{Un dato in più: la stima ricorsiva}
% ============================================================

Esplicitiamo che la stima è basata su $p$ dati:
\[ \vett{a}_{\text{LS}}(p) = \left( \matr{X}(p) \matr{X}^T(p) \right)^{-1} \matr{X}(p)\, \vett{y}(p) . \]
Consideriamo questo scenario: si vuole migliorare progressivamente la stima aggiungendo nuove osservazioni, man mano che
arrivano, e farlo con una complessità limitata. Supponiamo $p \geq n$ e di aver già calcolato $\vett{a}_{\text{LS}}(p)$.
Arrivano un nuovo vettore $\vett{x}(p + 1)$ e una nuova osservazione $\theta(p + 1)$; la nuova stima sarebbe
\[ \vett{a}_{\text{LS}}(p + 1) = \left( \matr{X}(p + 1) \matr{X}^T(p + 1) \right)^{-1} \matr{X}(p + 1)\, \vett{y}(p +
   1) . \]
Dobbiamo ricalcolare tutto da zero? No: la matrice da invertire cambia di poco, e la si può aggiornare.

\teorema[shermanmorrison]{Formula di Sherman-Morrison}{
  Sia $\matr{R}$ una matrice $n \times n$ invertibile e siano $\vett{u}$ e $\vett{v}$ vettori colonna $n$-dimensionali,
  con $1 + \vett{v}^T \matr{R}^{-1} \vett{u} \neq 0$. Allora l'inversa della matrice con \emph{aggiornamento di rango 1}
  $\matr{R} + \vett{u}\vett{v}^T$ è
  \[ \left( \matr{R} + \vett{u}\vett{v}^T \right)^{-1} = \matr{R}^{-1} - \frac{\matr{R}^{-1} \vett{u}\, \vett{v}^T
     \matr{R}^{-1}}{1 + \vett{v}^T \matr{R}^{-1} \vett{u}} . \]
}

\begin{dimostrazione}[basta moltiplicare]
\passo{Si moltiplica} Chiamiamo $c = \vett{v}^T \matr{R}^{-1} \vett{u}$: è un numero ($1 \times n$ per $n \times n$ per
$n \times 1$). Moltiplicando $\matr{R} + \vett{u}\vett{v}^T$ per il secondo membro:
\[ \left( \matr{R} + \vett{u}\vett{v}^T \right) \left( \matr{R}^{-1} - \frac{\matr{R}^{-1} \vett{u}\vett{v}^T
   \matr{R}^{-1}}{1 + c} \right) = \matr{I} + \vett{u}\vett{v}^T \matr{R}^{-1} - \frac{\vett{u}\vett{v}^T \matr{R}^{-1} +
   \vett{u} \left( \vett{v}^T \matr{R}^{-1} \vett{u} \right) \vett{v}^T \matr{R}^{-1}}{1 + c} . \]
\passo{Si raccoglie} Il numero $c$ si può spostare dove si vuole: il numeratore della frazione è $(1 + c)\,
\vett{u}\vett{v}^T \matr{R}^{-1}$, e la frazione vale $\vett{u}\vett{v}^T \matr{R}^{-1}$. Resta $\matr{I}$: il secondo
membro è proprio l'inversa.
\end{dimostrazione}

Le slide scrivono il denominatore come $1 + \vett{u}^T \matr{R}^{-1} \vett{v}$; nel nostro uso $\vett{u} = \vett{v}$ e
$\matr{R}$ è simmetrica, quindi le due scritture coincidono.

\teorema[rls]{La stima ai minimi quadrati ricorsiva}{
  Siano $\matr{R}(p) = \matr{X}(p)\matr{X}^T(p)$ e $K(p + 1) = \vett{x}(p + 1)^T \matr{R}^{-1}(p)\, \vett{x}(p + 1)$
  (un numero). Allora
  \[ \matr{R}^{-1}(p + 1) = \matr{R}^{-1}(p) - \frac{\matr{R}^{-1}(p)\, \vett{x}(p + 1)\, \vett{x}(p + 1)^T
     \matr{R}^{-1}(p)}{1 + K(p + 1)} , \]
  \[ \vett{a}(p + 1) = \left[ \matr{I}_n - \frac{\matr{R}^{-1}(p)\, \vett{x}(p + 1)\, \vett{x}(p + 1)^T}{1 + K(p + 1)}
     \right] \left[ \vett{a}(p) + \theta(p + 1)\, \matr{R}^{-1}(p)\, \vett{x}(p + 1) \right] . \]
  Il calcolo comporta $O(n^2)$ operazioni, indipendentemente da $p$: la complessità non cresce con il numero di dati.
}

\begin{dimostrazione}
\passo{La matrice come somma} Il prodotto $\matr{X}(p)\matr{X}^T(p)$, fatto «colonne per righe», è la somma dei
prodotti colonna per riga dei singoli vettori:
\[ \matr{X}(p)\matr{X}^T(p) = \left[ \vett{x}(1)\ \cdots\ \vett{x}(p) \right] \begin{bmatrix} \vett{x}^T(1) \\ \vdots
   \\ \vett{x}^T(p) \end{bmatrix} = \sum_{i=1}^{p} \vett{x}(i)\, \vett{x}^T(i) . \]
\passo{Un dato in più è un aggiornamento di rango 1} Aggiungendo il termine $i = p + 1$:
\[ \matr{R}(p + 1) = \sum_{i=1}^{p+1} \vett{x}(i)\vett{x}^T(i) = \matr{R}(p) + \underbrace{\vett{x}(p +
   1)}_{\vett{u}}\, \underbrace{\vett{x}^T(p + 1)}_{\vett{v}^T} , \]
e la formula di Sherman-Morrison dà l'inversa, con $\vett{u} = \vett{v} = \vett{x}(p + 1)$.
\passo{Il vettore dei termini noti} La matrice dei dati si allunga di una colonna e il vettore delle osservazioni di
un elemento: $\matr{X}(p + 1) = [\matr{X}(p)\ \ \vett{x}(p + 1)]$ e $\vett{y}(p + 1) = [\vett{y}^T(p),\ \theta(p +
1)]^T$. Moltiplicando a blocchi, $\matr{X}(p + 1)\vett{y}(p + 1) = \matr{X}(p)\vett{y}(p) + \theta(p + 1)\, \vett{x}(p +
1)$.
\passo{Si mette insieme} Scriviamo in breve $\vett{x}$ per $\vett{x}(p + 1)$, $\theta$ per $\theta(p + 1)$ e $K$ per
$K(p + 1)$. Raccogliendo $\matr{R}^{-1}(p)$ a destra nella formula dell'inversa (che si può fare perché
$\matr{R}^{-1}(p) = \matr{I}_n \matr{R}^{-1}(p)$, inserendo la matrice identità, che non cambia nulla):
\[ \vett{a}(p + 1) = \left[ \matr{I}_n - \frac{\matr{R}^{-1}(p)\, \vett{x}\vett{x}^T}{1 + K} \right] \matr{R}^{-1}(p)
   \left[ \matr{X}(p)\vett{y}(p) + \theta\, \vett{x} \right] , \]
e $\matr{R}^{-1}(p)\matr{X}(p)\vett{y}(p) = \vett{a}(p)$ è la stima precedente: si ottiene la formula dell'enunciato.
\passo{Il costo} Servono solo prodotti matrice per vettore e colonna per riga di dimensione $n$: ciascuno costa circa
$n^2$ moltiplicazioni, qualunque sia $p$. Rifare tutto da capo costerebbe invece circa $n^2 p$ operazioni per
$\matr{X}\matr{X}^T$ più circa $n^3$ per l'inversione.
\end{dimostrazione}

\info{La stessa formula, letta come una correzione}{
  Svolgendo il prodotto (con $\vett{x}$, $\theta$, $K$ in breve, $\matr{R}^{-1} = \matr{R}^{-1}(p)$, e ricordando che
  $\vett{x}^T \matr{R}^{-1} \vett{x} = K$ e $\vett{x}^T \vett{a}(p)$ sono numeri):
  \[ \vett{a}(p) + \theta \matr{R}^{-1}\vett{x} - \frac{\matr{R}^{-1}\vett{x} \left( \vett{x}^T \vett{a}(p) \right) + \theta
     \matr{R}^{-1}\vett{x}\, K}{1 + K} = \vett{a}(p) + \matr{R}^{-1}\vett{x} \left[ \theta \left( 1 - \frac{K}{1 + K}
     \right) - \frac{\vett{x}^T \vett{a}(p)}{1 + K} \right] , \]
  e siccome $1 - \frac{K}{1 + K} = \frac{1}{1 + K}$ la formula diventa
  \[ \vett{a}(p + 1) = \vett{a}(p) + \frac{\matr{R}^{-1}(p)\, \vett{x}(p + 1)}{1 + K(p + 1)} \Big[ \theta(p + 1) -
     \vett{x}(p + 1)^T \vett{a}(p) \Big] . \]
  Il termine fra parentesi quadre è l'errore che il vecchio modello commette sul dato nuovo (l'\emph{innovazione}):
  se il modello aveva già previsto bene $\theta(p + 1)$, la stima quasi non cambia; se aveva sbagliato, viene corretta
  in proporzione all'errore. È la stessa idea del filtro di Kalman.
}

\begin{esempio}{La stima che migliora dato dopo dato}
Si stimano i coefficienti della retta $\theta = a_1 + a_2 t$ da dati simulati con $a_1 = 1$, $a_2 = 2$ e un rumore
gaussiano di deviazione standard 1, che arrivano uno alla volta. La prima stima si calcola con la formula diretta sui
primi due dati (servono $p \geq n = 2$); poi si usa la formula ricorsiva.
\begin{center}
\begin{tikzpicture}
\begin{axis}[width=12cm, height=5.6cm, xmin=0, xmax=61, ymin=0, ymax=3.8, xtick={2,10,20,30,40,50,60},
  ytick={0,1,2,3}, xlabel={numero di dati $p$}, axis lines=left, tick label style={font=\scriptsize},
  label style={font=\small}, grid=major, grid style={mspigray!20}, legend style={font=\scriptsize, at={(0.99,0.98)},
  anchor=north east}, legend columns=2]
  \addplot[mspiA, line width=1.1pt, mark=*, mark size=1pt] coordinates {(2,3.522) (3,1.011) (4,0.995) (5,0.758) (6,0.764) (7,1.402) (8,0.429) (9,0.647) (10,0.613) (11,0.748) (12,0.751) (13,0.704) (14,0.868) (15,0.788) (16,0.613) (17,0.684) (18,0.689) (19,0.689) (20,0.889) (21,0.895) (22,0.793) (23,0.727) (24,0.757) (25,0.691) (26,0.435) (27,0.510) (28,0.511) (29,0.491) (30,0.487) (31,0.479) (32,0.552) (33,0.572) (34,0.541) (35,0.513) (36,0.597) (37,0.657) (38,0.691) (39,0.696) (40,0.718) (41,0.713) (42,0.699) (43,0.683) (44,0.717) (45,0.667) (46,0.622) (47,0.622) (48,0.624) (49,0.634) (50,0.613) (51,0.556) (52,0.620) (53,0.622) (54,0.628) (55,0.685) (56,0.658) (57,0.653) (58,0.730) (59,0.758) (60,0.770)};
  \addplot[mspiB, line width=1.1pt, mark=*, mark size=1pt] coordinates {(2,1.387) (3,1.954) (4,1.942) (5,1.999) (6,1.924) (7,1.821) (8,2.075) (9,2.075) (10,2.079) (11,2.049) (12,2.046) (13,2.079) (14,2.041) (15,2.036) (16,2.137) (17,2.120) (18,2.119) (19,2.129) (20,2.016) (21,2.012) (22,2.074) (23,2.092) (24,2.084) (25,2.090) (26,2.157) (27,2.155) (28,2.155) (29,2.160) (30,2.154) (31,2.140) (32,2.128) (33,2.127) (34,2.135) (35,2.126) (36,2.105) (37,2.094) (38,2.086) (39,2.086) (40,2.073) (41,2.080) (42,2.075) (43,2.078) (44,2.070) (45,2.082) (46,2.116) (47,2.122) (48,2.122) (49,2.108) (50,2.121) (51,2.135) (52,2.120) (53,2.118) (54,2.117) (55,2.100) (56,2.098) (57,2.101) (58,2.079) (59,2.071) (60,2.064)};
  \addplot[mspiA, dashed, domain=0:61, samples=2] {1};
  \addplot[mspiB, dashed, domain=0:61, samples=2] {2};
  \legend{{$a_1(p)$}, {$a_2(p)$}, {valore vero 1}, {valore vero 2}}
\end{axis}
\end{tikzpicture}
\end{center}
Con due soli dati la stima è molto lontana dal vero ($a_1 \approx 3{,}5$, $a_2 \approx 1{,}4$), ma bastano pochi dati in
più per avvicinarsi. Dopo 60 dati la stima ricorsiva coincide, fino all'ultima cifra, con quella calcolata da capo su
tutti i dati.
\end{esempio}

\gbox{In sintesi}{mspiA!80!black}{
\begin{itemize}
  \item Statistica descrittiva: insieme di addestramento $\matr{X} = [\vett{x}(1) \cdots \vett{x}(p)]$ e osservazioni
        $\vett{y}$; modello lineare $\theta(i) = \vett{a}^T\vett{x}(i) + \epsilon_i$.
  \item \textbf{Minimi quadrati}: si minimizza $\norm{\matr{X}^T\vett{a} - \vett{y}}^2$; $\vett{a}_{\text{LS}} =
        (\matr{X}\matr{X}^T)^{-1}\matr{X}\vett{y}$, se $p \geq n$ e le caratteristiche sono linearmente indipendenti.
        È lo stimatore lineare MMSE con le medie campione.
  \item \textbf{Sherman-Morrison}: $(\matr{R} + \vett{u}\vett{v}^T)^{-1} = \matr{R}^{-1} - \frac{\matr{R}^{-1}\vett{u}
        \vett{v}^T\matr{R}^{-1}}{1 + \vett{v}^T\matr{R}^{-1}\vett{u}}$.
  \item \textbf{Stima ricorsiva}: $\matr{R}(p + 1) = \matr{R}(p) + \vett{x}(p + 1)\vett{x}^T(p + 1)$; l'aggiornamento
        di $\matr{R}^{-1}$ e di $\vett{a}$ costa $O(n^2)$ operazioni, indipendentemente da $p$.
\end{itemize}
}
